\documentclass[12pt]{article}
\usepackage{amsmath, amsthm, amssymb}
\usepackage[utf8]{inputenc}
\usepackage[margin=1in]{geometry}

\newtheorem{theorem}{Theorem}

\begin{document}

\section{Definable Completeness}
\begin{theorem}
Let $\mathcal{M}$ be an o-minimal structure. Every non-empty definable set $S \subseteq M$ that is bounded above has a supremum in $M$.
\end{theorem}
\begin{proof}
Because $S$ is definable, o-minimality dictates it can be written exactly as a finite union of open intervals and discrete points:
\[ S = I_1 \cup \dots \cup I_k \cup P \]
where $P$ is a finite set.

Because $S$ is bounded above, every component of $S$ must also be bounded above. 
\begin{enumerate}
    \item For the finite set $P$, by Theorem 1, it has a maximum, which serves as its supremum: $\sup(P) = \max(P)$.
    \item For any open interval $I_j$, because it is bounded above, it must be of the form $(a_j, b_j)$ or $(-\infty, b_j)$ for some endpoint $b_j \in M$. The supremum of this interval is exactly its right endpoint, $b_j$.
\end{enumerate}

Therefore, every individual component of $S$ has a supremum in $M$. By Theorem 3, the finite union of sets that each have a supremum also has a supremum. Thus, $S$ has a supremum in $M$.
\end{proof}

\end{document}